% Fonte principal autossuficiente. Compilar com pdfLaTeX, XeLaTeX ou Tectonic.
% As quebras de página organizam os argumentos dentro do limite de quinze páginas A4.
\documentclass[11pt,a4paper]{article}
\usepackage[T1]{fontenc}
\usepackage{iftex}
\ifPDFTeX\usepackage[utf8]{inputenc}\fi
\usepackage[brazil]{babel}
\addto\captionsbrazil{\renewcommand{\refname}{Bibliografia}}
\usepackage{lmodern}
\usepackage[margin=2.15cm,headheight=14pt,headsep=0.55cm,footskip=0.8cm]{geometry}
\usepackage{amsmath,amssymb,amsthm,mathtools}
\usepackage{microtype,booktabs,array,enumitem}
\usepackage{xcolor,tikz}
\usetikzlibrary{arrows.meta,positioning,calc}
\usepackage{fancyhdr,caption}
\usepackage[unicode,colorlinks=true,linkcolor=blue!45!black,citecolor=blue!45!black,urlcolor=blue!45!black]{hyperref}
\hypersetup{pdftitle={Árvores e grafos de estados em uma variante competitiva de Hi-Lo},pdfauthor={Caio Ladislau do Couto; João Henrique Machado; Ricardo Rodrigues Tose; Miguel Vitor Pereira}}
\urlstyle{same}
\setlength{\parindent}{1.1em}
\setlength{\parskip}{3pt}
\setlength{\emergencystretch}{1.5em}
\setlist{nosep,leftmargin=1.6em,topsep=4pt}
\captionsetup{font=small,labelfont=bf,skip=5pt,hypcap=false}
\pagestyle{fancy}
\fancyhf{}
\fancyhead[L]{\small Hi-Lo competitivo: árvores e estados}
\fancyhead[R]{\small Matemática Discreta}
\fancyfoot[C]{\thepage}
\renewcommand{\headrulewidth}{0.3pt}
\newtheorem{proposicao}{Proposição}[section]
\newtheorem{lema}[proposicao]{Lema}
\newtheorem{teorema}[proposicao]{Teorema}
\newtheorem{corolario}[proposicao]{Corolário}
\theoremstyle{definition}
\newtheorem{definicao}{Definição}[section]
\newtheorem{observacao}[definicao]{Observação}
\newcommand{\grauout}{\deg^{+}}
\newcommand{\grauin}{\deg^{-}}
\newcommand{\A}{\mathrm{A}}
\newcommand{\Er}{\mathrm{E}}
\newcommand{\destaque}[1]{\par\smallskip\noindent\textbf{#1}\enspace}
\newcommand{\proxpagina}{\newpage}
\raggedbottom

\begin{document}
% PÁGINA 1 ---------------------------------------------------------------
\thispagestyle{plain}
\begin{center}
{\LARGE\bfseries Árvores e grafos de estados\\[3pt]
em uma variante competitiva de Hi-Lo\par}
\vspace{0.45cm}
Caio Ladislau do Couto \quad João Henrique Machado\\
Ricardo Rodrigues Tose \quad Miguel Vitor Pereira\\[6pt]
\textbf{Professora:} Maria Soledad Aronna\\
Matemática Discreta \quad FGV \quad Novembro de 2026
\end{center}
\vspace{0.15cm}
\begin{abstract}
Estudamos uma variante de Hi-Lo na qual dois jogadores preveem a relação entre a próxima carta e a referência atual. Um acerto pontua e conserva a vez; um erro transfere a vez. Representamos as partidas por uma árvore de históricos e reunimos históricos equivalentes em um grafo de estados. Demonstramos que esse grafo é dirigido acíclico e descrevemos, por um invariante de paridade, todos os pares de diferença de pontos e jogador da vez em cada camada. Essa descrição fornece uma fórmula fechada para o número de estados de qualquer composição inicial de cartas. Contamos também as folhas da árvore e os estados terminais. Provamos que o grafo associado é bipartido, é uma árvore exatamente quando restam no máximo duas cartas e, nos demais casos, contém um ciclo mínimo de quatro arestas. Relacionamos entradas múltiplas, árvores geradoras e arestas excedentes. Os resultados são ilustrados em finais reduzidos e aplicados ao baralho original de 36 cartas, cujo grafo, após a referência inicial, possui 402.187.501 estados. Essa contagem é obtida pela fórmula, sem construir o grafo completo.
\end{abstract}
\section{Introdução}
Uma partida contém tanto uma sequência de acontecimentos quanto uma situação atual. Duas sequências diferentes podem deixar as mesmas cartas disponíveis, a mesma referência, a mesma diferença de pontos e o mesmo jogador da vez. Quando isso ocorre, as possibilidades de continuação coincidem. Podemos, portanto, estudar duas representações: uma árvore que distingue os históricos e um grafo que identifica os estados iguais.

Escolhemos um jogo de cartas porque suas regras permitem definir os grafos por um procedimento explícito e, ao mesmo tempo, demonstrar resultados gerais. As perguntas centrais são: \emph{quantos estados são alcançáveis, como se compara esse número ao de históricos e quando a reunião de históricos produz ciclos no grafo associado?} Os exemplos desenvolvidos no texto mostram como as definições e as fórmulas respondem a essas perguntas em composições concretas.

A parte prática do projeto é apresentada em um arquivo HTML separado, dedicado à visualização de exemplos dos grafos estudados neste relatório. As definições, demonstrações e exemplos aqui desenvolvidos fornecem a fundamentação matemática dessa visualização.

A Seção~2 fixa regras, conceitos e notação. A Seção~3 apresenta os sete resultados principais e suas demonstrações. A Seção~4 desenvolve exemplos e aplica as fórmulas ao baralho completo. A Seção~5 discute o alcance dos resultados. O recorte é estrutural: contar possibilidades e estudar conexões não exige determinar uma estratégia ótima nem atribuir probabilidades aos estados.

\proxpagina
% PÁGINA 2 ---------------------------------------------------------------
\section{Propriedades e definições básicas}
\subsection{Regras da variante competitiva}
O baralho original contém 36 cartas: quatro cópias de cada valor do conjunto
\[
R=\{2,3,\ldots,10\}.
\]
Os naipes e as identidades físicas das cópias não interferem no jogo. O embaralhamento é uniforme e não há reposição. Os jogadores conhecem a composição inicial e todas as cartas já reveladas, mas desconhecem a ordem das cartas ainda ocultas. As regras são:
\begin{enumerate}
\item Ambos começam com zero ponto. Revela-se uma carta inicial, que passa a ser a \emph{referência}, sem atribuição de ponto. Designa-se o primeiro jogador.
\item Antes de revelar outra carta, o jogador da vez escolhe uma previsão: \emph{menor}, \emph{igual} ou \emph{maior}, sempre em relação à referência atual.
\item Revela-se a próxima carta. Se a previsão estiver correta, o jogador recebe um ponto e mantém a vez. Caso contrário, ninguém recebe ponto e a vez passa ao adversário. Um erro não retira pontos anteriores.
\item Em ambos os casos, a carta revelada torna-se a nova referência e deixa o estoque oculto. Repete-se o procedimento enquanto houver cartas.
\item A última previsão também é avaliada. Em seguida, comparam-se as pontuações: vence quem tiver mais pontos, admitindo-se empate.
\end{enumerate}

\destaque{Uma partida curta.}
Considere referência 6, cartas restantes $[6,7,7]$ e jogador 1 começando. Fixemos, para ilustrar, a ordem oculta $7,7,6$. Os jogadores não conhecem essa ordem quando escolhem suas previsões.
\begin{center}
\small
\begin{tabular}{@{}ccclccc@{}}
\toprule
Jogada & Jogador & Referência & Previsão & Carta & Placar após & Próxima vez\\
\midrule
1 & 1 & 6 & maior & 7 & $1\!:\!0$ & 1\\
2 & 1 & 7 & maior & 7 & $1\!:\!0$ & 2\\
3 & 2 & 7 & menor & 6 & $1\!:\!1$ & 2\\
\bottomrule
\end{tabular}
\captionof{table}{A igualdade na segunda jogada produz erro. A partida termina empatada.}
\end{center}

\destaque{Hipóteses essenciais às contagens.}
Todas as três previsões são sempre permitidas, inclusive as que certamente erram. Por exemplo, com referência 10, a previsão \emph{maior} continua permitida. Estudamos todas as partidas compatíveis com as regras, sem restringir os jogadores a escolhas consideradas vantajosas. Uma carta, por sua vez, só pode ser revelada se ainda existir no estoque.

A referência inicial já foi retirada \emph{antes} da primeira jogada avaliada. Assim, um baralho de 36 cartas produz 35 jogadas. Nos finais reduzidos, o número $N$ sempre designará as cartas \emph{restantes na situação inicial escolhida}, e não o tamanho do baralho anterior à referência. A raiz corresponde a essa situação, antes de qualquer nova previsão.

\proxpagina
% PÁGINA 3 ---------------------------------------------------------------
\subsection{Vocabulário de grafos}
Adotamos a terminologia de caminhos das anotações da disciplina~\cite{nina}. As notas de Rangel, Oliveira e Araujo~\cite{rangel} oferecem a referência complementar para digrafos, bipartição e árvores geradoras. Todos os grafos considerados aqui são finitos.

\begin{definicao}[Grafos e orientações]
Um \emph{grafo simples} $G=(V,E)$ tem conjunto de vértices $V$ e conjunto de arestas $E$ formado por pares não ordenados $\{u,v\}$ com $u\ne v$. Assim, não há laços nem arestas paralelas. Em um \emph{grafo dirigido}, $E$ é um conjunto de pares ordenados $(u,v)$, escritos $u\to v$. O \emph{grafo associado} é obtido ao esquecer as orientações.
\end{definicao}

\begin{definicao}[Caminhos, ciclos e conexidade]
Um caminho de comprimento $\ell$ é uma sequência alternada $(u_0,e_1,u_1,\ldots,e_\ell,u_\ell)$ de vértices e arestas incidentes, admitindo repetições; para $\ell=0$, há apenas um vértice. Um caminho é \emph{simples} se não repete vértices. Um ciclo é um caminho fechado de comprimento positivo sem repetição de arestas; é \emph{simples} se só repete o vértice inicial ao final. Um grafo é \emph{conexo} se qualquer par de vértices pode ser ligado por um caminho.
\end{definicao}

Em grafos dirigidos, os caminhos respeitam o sentido das arestas. Um digrafo é \emph{acíclico} quando não possui ciclos dirigidos. Os graus $\grauin(v)$ e $\grauout(v)$ contam, respectivamente, as arestas que chegam a $v$ e as que saem de $v$.

\begin{observacao}[Convenção terminológica]
O que chamamos de caminho com possíveis repetições é denominado \emph{passeio} em diversas referências, inclusive nas notas complementares. A distinção entre caminho e caminho simples será mantida explicitamente neste texto.
\end{observacao}

\begin{definicao}[Árvores e bipartição]
Uma \emph{árvore} é um grafo não vazio no qual existe um único caminho simples entre quaisquer dois vértices. Uma árvore \emph{com raiz} tem um vértice designado como raiz. O pai de um vértice não inicial é seu vizinho imediatamente anterior no caminho vindo da raiz; os sucessores imediatos são seus filhos. Um vértice sem filhos é uma \emph{folha}. Uma \emph{árvore geradora} de um grafo conexo é um subgrafo que contém todos os seus vértices e é uma árvore. Um grafo é \emph{bipartido} se $V$ pode ser separado em dois conjuntos disjuntos e toda aresta tem um extremo em cada conjunto.
\end{definicao}

Usaremos a equivalência elementar entre ser árvore e ser conexo sem ciclos simples. De fato, em um grafo conexo existe um caminho simples entre cada par: basta eliminar trechos entre repetições. Se houvesse dois caminhos simples distintos entre o mesmo par, um trecho de separação seguido de reencontro produziria um ciclo simples. Reciprocamente, um ciclo fornece dois caminhos simples entre dois de seus vértices.

No nosso modelo, a existência de ciclos no grafo associado não contradirá a aciclicidade dirigida: percorrer um ciclo depois de apagar as setas não significa poder realizar essa sequência de jogadas.

\proxpagina
% PÁGINA 4 ---------------------------------------------------------------
\subsection{Estados, históricos e notação}
\begin{definicao}[Estado do jogo]
Um estado é a tupla
\[
s=(c,r,d,t),\qquad c=(c_2,c_3,\ldots,c_{10}),
\]
na qual $c_j$ é a quantidade restante de cartas de valor $j$, $r$ é a referência atual, $d$ é a pontuação do jogador 1 menos a do jogador 2, e $t\in\{1,2\}$ indica quem joga. Escrevemos $n(s)=\sum_{j=2}^{10}c_j$. O estado é \emph{terminal} quando $n(s)=0$.
\end{definicao}

O vencedor é determinado pelo sinal final de $d$. Guardar a diferença basta para atualizar o resultado relativo: um acerto de J1 aumenta $d$ em um; um acerto de J2 o diminui em um; um erro o conserva. Não é necessário guardar as duas pontuações separadamente. Em um terminal, mantemos $r$ e $t$ na tupla por convenção, embora já não haja próxima jogada.

Um \emph{histórico} registra, em ordem, cada previsão e cada valor revelado desde a raiz. Cópias do mesmo valor não são distinguidas. Na \emph{árvore de históricos}, cada histórico é um vértice, o histórico vazio é a raiz, e cada aresta acrescenta uma jogada. No \emph{grafo de estados}, cada estado alcançável aparece uma única vez e $s\to u$ é uma aresta se alguma jogada leva de $s$ a $u$. Havendo várias previsões com o mesmo sucessor, preservamos seus rótulos em uma única aresta.

Por exemplo, da raiz com referência 6 e cartas $[6,7,7]$, prever \emph{menor} e revelar 7, ou prever \emph{igual} e revelar 7, são históricos diferentes. Ambos deixam $[6,7]$, referência 7, $d=0$ e $t=2$. Há dois vértices na árvore, mas um só estado no grafo. Esse estado não guarda os dois históricos dentro de si: guarda as informações comuns que bastam à continuação.

\destaque{Notação para contar estados.}
Sejam $v_1,\ldots,v_k$ os valores distintos presentes no estoque da raiz. Adotamos a seguinte notação:
\begin{center}
\begin{tabular}{@{}cl@{}}
\toprule
Símbolo & Significado\\
\midrule
$C_i$ & Quantidade de cartas de valor $v_i$ no estoque da raiz.\\
$N=\sum_i C_i$ & Número de cartas restantes na raiz, após a referência.\\
$a=(a_1,\ldots,a_k)$ & Vetor das quantidades consumidas desde essa raiz.\\
$|a|=\sum_i a_i=m$ & Número de jogadas realizadas; camada atual.\\
$c_{v_i}=C_i-a_i$ & Quantidade que ainda resta do valor $v_i$.\\
$V_m$ & Conjunto dos estados alcançáveis após $m$ jogadas.\\
\bottomrule
\end{tabular}
\end{center}
Os índices $i$ numeram \emph{categorias}; $v_i$ é o \emph{valor da carta}. As coordenadas de $a$ têm posições fixas, mas não registram a ordem temporal das revelações. O símbolo $|a|$ indica soma de coordenadas; em $|V_m|$, as barras indicam cardinalidade. As categorias ausentes do estoque são omitidas, de modo que $C_i>0$.

\proxpagina
% PÁGINA 5 ---------------------------------------------------------------
\section{Resultados principais}
Fixamos uma raiz com diferença $d=0$, jogador 1 começando e $N$ cartas restantes. Uma diferença inicial distinta apenas translada $d$; se começa J2, a troca dos nomes dos jogadores reduz a essa convenção. As contagens e as propriedades estruturais permanecem as mesmas. A fórmula de paridade abaixo é escrita especificamente nessa convenção.

\subsection{Representações, camadas e graus}
\begin{proposicao}[Representações, camadas e graus]\label{prop:representacoes}
Os históricos formam uma árvore com raiz. Históricos com o mesmo estado admitem as mesmas continuações e podem ser reunidos no grafo de estados. Nesse grafo:
\begin{enumerate}[label=\roman*)]
\item cada aresta liga a camada $m$ à camada $m+1$;
\item todo caminho dirigido é simples, e os caminhos da raiz que não podem ser prolongados têm comprimento $N$;
\item havendo $q$ valores distintos disponíveis, o grau de saída é $2q$, enquanto cada histórico correspondente tem $3q$ filhos na árvore.
\end{enumerate}
O grafo dirigido é acíclico e seu associado é conexo.
\end{proposicao}

\begin{proof}
\textbf{A árvore.} Apagar a última jogada de um histórico não vazio produz seu único pai. Repetindo essa operação chegamos ao histórico vazio, logo a representação é conexa. Se existisse um ciclo simples, um histórico de comprimento máximo nesse ciclo teria dois vizinhos de comprimento uma unidade menor: seriam dois pais distintos. Portanto, não há ciclos simples e temos uma árvore, pela equivalência da Seção~2.2.

\textbf{A reunião dos históricos.} Se dois históricos terminam em $(c,r,d,t)$, têm os mesmos valores disponíveis. Fixadas uma previsão e uma próxima carta possível, a comparação com $r$ produz o mesmo acerto ou erro. As regras atualizam igualmente as quatro componentes, dando o mesmo sucessor. Por indução no comprimento da continuação, as mesmas sequências futuras são possíveis e produzem os mesmos estados. Essa é a justificativa para a identificação.

\textbf{As camadas.} Depois de $m$ jogadas, $n(s)=N-m$. Uma jogada consome exatamente uma carta; assim, toda aresta avança uma camada. Um caminho dirigido não pode repetir um vértice, pois não pode recuperar cartas. Logo, não existe ciclo dirigido. Enquanto $n(s)>0$, há carta possível e alguma previsão permitida; somente em $m=N$ um caminho vindo da raiz deixa de admitir prolongamento. Todo estado é alcançável da raiz por definição; apagando as setas, todos ficam conectados a ela.

\textbf{Os graus.} Para cada um dos $q$ valores reveláveis há três previsões, dando $3q$ filhos na árvore. Exatamente uma acerta e duas erram. As duas erradas levam ao mesmo sucessor; acerto e erro levam a sucessores distintos, pois a vez é, respectivamente, conservada e trocada. Valores revelados distintos dão referências distintas. Há, portanto, exatamente $2q$ sucessores no grafo de estados. Para um terminal, $q=0$.
\end{proof}

Ao apagar as orientações, não se fundem arestas distintas: não existem setas em sentidos opostos entre os mesmos vértices, pois todas avançam uma camada. Logo, o grafo dirigido e seu associado têm as mesmas quantidades $|V|$ e $|E|$.

\proxpagina
% PÁGINA 6 ---------------------------------------------------------------
\subsection{Paridade e pares de placar e jogador}
Escrevemos $\A$ para acerto e $\Er$ para erro. São \emph{resultados}, não previsões. Fixada uma ordem viável das cartas, qualquer palavra de $\A$ e $\Er$ pode ocorrer em algum histórico: em cada posição há uma previsão correta e duas incorretas. Isso prova existência de históricos; não fornece uma estratégia capaz de conhecer a próxima carta.

\begin{lema}[Pares de placar e jogador]\label{lem:pares}
Após $1\le m\le N$ jogadas, para qualquer ordem viável dos valores revelados, existem exatamente $2m$ pares $(d,t)$ possíveis. As diferenças percorrem todos os inteiros de $1-m$ até $m$, e
\[
t=1\quad\Longleftrightarrow\quad m-d\text{ é par}.
\]
\end{lema}

\begin{proof}
\textbf{1. Restrição de paridade.} Na raiz, $m-d=0$ e $t=1$. Verifiquemos os efeitos de uma jogada:
\begin{center}
\begin{tabular}{@{}lcc@{}}
\toprule
Resultado & Novo valor de $m-d$ & Vez\\
\midrule
J1 acerta & $(m+1)-(d+1)=m-d$ & Conserva\\
J2 acerta & $(m+1)-(d-1)=m-d+2$ & Conserva\\
Qualquer jogador erra & $(m+1)-d=m-d+1$ & Troca\\
\bottomrule
\end{tabular}
\end{center}
A paridade muda exatamente quando a vez muda. Por indução nas jogadas, a equivalência é preservada. Assim, para cada diferença $d$, há no máximo um jogador da vez possível.

\textbf{2. Limites da diferença.} J1 marca no máximo $m$ pontos, logo $d\le m$. Para J2 chegar a pontuar, J1 precisa ter errado antes; essa jogada não pontua. Portanto, J2 marca no máximo $m-1$ pontos, dando $d\ge -(m-1)=1-m$. Há $m-(1-m)+1=2m$ inteiros nesse intervalo. Junto à paridade, isso prova o limite superior de $2m$ pares.

\textbf{3. Realização de todos os pares.} Admitimos expoentes zero nas palavras abaixo. Nas palavras $\A^a\Er\A^b$, com $a+b=m-1$, J1 marca $a$ pontos, erra e passa a vez a J2, que marca $b$ pontos. Terminamos com J2 na vez e diferença $d=a-b=2a-(m-1)$. Para $a=0,\ldots,m-1$, obtemos $m$ diferenças distintas:
\[
1-m,\ 3-m,\ldots,\ m-1.
\]
Para $m\ge2$, as palavras $\A^a\Er\A^b\Er$, com $a+b=m-2$, acrescentam um erro de J2, devolvendo a vez a J1 sem modificar o placar. Agora $d=a-b=2a-(m-2)$, e $a=0,\ldots,m-2$ realiza $m-1$ diferenças distintas:
\[
2-m,\ 4-m,\ldots,\ m-2.
\]
Falta o par $(m,1)$, realizado por $\A^m$. As famílias fornecem $m+(m-1)+1=2m$ pares distintos. Para $m=1$, $\A$ e $\Er$ dão $(1,1)$ e $(0,2)$. Todas essas palavras são compatíveis com a ordem de cartas fixada, pela observação anterior. O limite superior é atingido.
\end{proof}

Fixada a raiz, $m=N-n(s)$ é recuperado do estoque. O lema mostra que $t$ também pode ser recuperado de $m$ e $d$. Mantemos $t$ no estado para tornar explícitas as regras de pontuação e troca de vez; essa redundância não altera a contagem.

\proxpagina
% PÁGINA 7 ---------------------------------------------------------------
\subsection{Contagem exata dos estados}
O lema não conta estados completos: um mesmo par $(d,t)$ pode acompanhar estoques ou referências diferentes. Para contar os estados, precisamos combinar esse par com as informações das cartas.

As provas valem para qualquer multiconjunto finito de valores e multiplicidades positivas; não utilizam o limite de quatro cópias do baralho original. A referência inicial pode estar ausente do estoque restante.

Por exemplo, para $m=3$ o lema fornece
\[
(-2,2),\ (-1,1),\ (0,2),\ (1,1),\ (2,2),\ (3,1).
\]
Fixar também um consumo e uma referência produzirá seis estados completos distintos. Já trocar o consumo ou a referência produzirá outros estados, ainda que esses mesmos seis pares reapareçam.

\begin{teorema}[Número de vértices]\label{teo:vertices}
Se o estoque da raiz possui $k$ valores distintos, de multiplicidades positivas $C_1,\ldots,C_k$, e $N=\sum_{i=1}^{k}C_i\ge1$, então
\begin{equation}\label{eq:vertices}
\boxed{\displaystyle |V|=1+(N+1)\sum_{j=1}^{k} C_j\prod_{i\ne j}(C_i+1).}
\end{equation}
O número depende apenas das multiplicidades, sem depender dos valores numéricos das cartas ou da referência inicial. Um produto sem fatores vale 1.
\end{teorema}

\begin{proof}
Fixemos uma referência atual $v_j$ após ao menos uma jogada. Ela precisa ter sido revelada; por isso, seu consumo satisfaz $a_j\ge1$. As demais coordenadas satisfazem apenas os limites do estoque. Definimos
\[
B_j=\bigl\{a\in\mathbb Z^k:1\le a_j\le C_j,\quad
0\le a_i\le C_i\ (i\ne j)\bigr\}.
\]
Essas condições são também suficientes: para qualquer $a\in B_j$, ordenamos as cartas consumidas colocando uma cópia de $v_j$ por último. Assim, existe uma ordem viável de $|a|$ revelações com esse consumo e essa referência.

Fixados $a$ e $j$, temos $c_{v_i}=C_i-a_i$ e $r=v_j$. O Lema~\ref{lem:pares} fornece exatamente $2|a|$ pares $(d,t)$ para essa ordem, todos realizáveis. Como qualquer histórico com esses mesmos consumo e referência obedece ao mesmo lema, há exatamente $2|a|$ estados completos com essas informações.

Não há duplicação ao variar $a$ e $j$: vetores de consumo distintos deixam estoques distintos; índices $j$ distintos dão referências distintas. A raiz, por sua vez, tem consumo zero e não está em nenhum $B_j$.

Pelo princípio multiplicativo, o número de vetores em $B_j$ é
\[
|B_j|=C_j\prod_{i\ne j}(C_i+1).
\]
A coordenada $a_j$ tem $C_j$ escolhas; cada outra tem $C_i+1$, pois inclui o zero. Essa cardinalidade conta \emph{vetores de consumo}, e não estados. Para obter os estados com referência $v_j$, precisamos calcular
\[
S_j=\sum_{a\in B_j}2|a|.
\]
\proxpagina
% PÁGINA 8 ---------------------------------------------------------------
\destaque{Continuação: avaliação da soma.}
O peso $2|a|$ varia ao longo da caixa $B_j$. Para somá-lo, emparelhamos consumos por uma simetria da caixa. Definimos $a\mapsto a'$ coordenada a coordenada:
\begin{equation}\label{eq:espelho}
a'_j=C_j+1-a_j,\qquad a'_i=C_i-a_i\quad(i\ne j).
\end{equation}
Se $1\le a_j\le C_j$, então $1\le a'_j\le C_j$. Para as demais coordenadas, $0\le a_i\le C_i$ implica $0\le a'_i\le C_i$. Portanto, $a'\in B_j$. Além disso,
\[
C_j+1-(C_j+1-a_j)=a_j,\qquad C_i-(C_i-a_i)=a_i.
\]
Aplicar a transformação duas vezes devolve o vetor original. Logo, ela é uma bijeção de $B_j$ em si mesma, cuja inversa é ela própria. Somando as coordenadas de~\eqref{eq:espelho}, obtemos
\begin{equation}\label{eq:somasespelho}
|a'|=\sum_{i=1}^{k}C_i+1-\sum_{i=1}^{k}a_i=N+1-|a|.
\end{equation}
Como a bijeção percorre os mesmos vetores, apenas em outra ordem,
\begin{align*}
S_j&=\sum_{a\in B_j}2|a'|
=\sum_{a\in B_j}2(N+1-|a|)\\
&=2(N+1)|B_j|-S_j.
\end{align*}
Assim, $2S_j=2(N+1)|B_j|$, ou seja,
\[
S_j=(N+1)|B_j|.
\]
Somamos os grupos de referências, que são disjuntos, e acrescentamos a raiz:
\[
|V|=1+\sum_{j=1}^{k}S_j
=1+(N+1)\sum_{j=1}^{k}C_j\prod_{i\ne j}(C_i+1).
\]
Isso conclui a demonstração.
\end{proof}

\destaque{Exemplo do emparelhamento.}
Para $C=(1,2)$ e referência $v_2$, a caixa $B_2$ e seu espelhamento são:
\begin{center}
\begin{tabular}{@{}ccccc@{}}
\toprule
$a$ & $|a|$ & $a'$ & $|a'|$ & $|a|+|a'|$\\
\midrule
$(0,1)$ & 1 & $(1,2)$ & 3 & 4\\
$(0,2)$ & 2 & $(1,1)$ & 2 & 4\\
$(1,1)$ & 2 & $(0,2)$ & 2 & 4\\
$(1,2)$ & 3 & $(0,1)$ & 1 & 4\\
\bottomrule
\end{tabular}
\end{center}
Aqui $N=3$, $|B_2|=4$ e $S_2=2(1+2+2+3)=16=(N+1)|B_2|$. Não se afirma que $|a|=|a'|$ para cada vetor. A relação individual é $|a|+|a'|=N+1$; a bijeção preserva a \emph{soma sobre toda a caixa}. O argumento continua válido se algum vetor for seu próprio espelho, pois não exige dividir a caixa em pares distintos.

A referência \emph{atual} é uma das informações contadas, mas o valor da referência \emph{inicial} não aparece na fórmula. Após uma jogada, a referência é sempre um valor consumido; para cada ordem, todos os padrões de acerto e erro continuam realizáveis. Por isso, alterar a referência inicial muda quais previsões acertam, mas não a quantidade de estados alcançáveis.

\proxpagina
% PÁGINA 9 ---------------------------------------------------------------
\subsection{Contagem por camada e estados terminais}
\begin{corolario}[Camadas, terminais e cartas iguais]\label{cor:camadas}
A camada zero possui $|V_0|=1$. Para $1\le m\le N$,
\begin{equation}\label{eq:camadas}
\boxed{\displaystyle |V_m|=2m\!\sum_{\substack{0\le a_i\le C_i\ (1\le i\le k)\\a_1+\cdots+a_k=m}}\!|\{j:a_j>0\}|.}
\end{equation}
Há exatamente $2Nk$ estados terminais. Se todas as $N\ge1$ cartas restantes têm o mesmo valor, então $|V|=1+N(N+1)$.
\end{corolario}

\begin{proof}
Fixemos $m$ e um consumo $a$ de soma $m$. A referência atual precisa ser um valor consumido. Reciprocamente, qualquer valor com $a_j>0$ pode ser colocado na última posição de uma ordem dessas cartas. Assim, o conjunto $\{j:a_j>0\}$ tem exatamente os índices das referências possíveis para esse consumo.

Cada escolha de referência admite $2m$ pares $(d,t)$ pelo Lema~\ref{lem:pares}. Para esse $a$, há, portanto, $2m|\{j:a_j>0\}|$ estados. Consumos diferentes deixam estoques diferentes, logo os conjuntos contados são disjuntos. Somar todos os consumos viáveis fornece~\eqref{eq:camadas}. O fator $2m$ pode sair da soma porque $m$ é fixo.

No final, $m=N$ e o único consumo é $a=(C_1,\ldots,C_k)$. Todos os $k$ índices têm consumo positivo, de modo que $|V_N|=2Nk$. Se $k=1$, para cada $m\ge1$ existe apenas o consumo $a_1=m$ e uma referência possível. Logo, $|V_m|=2m$ e
\[
|V|=1+\sum_{m=1}^{N}2m=1+2\frac{N(N+1)}2=1+N(N+1).
\]
\end{proof}

\destaque{Como ler o somatório.}
Ele percorre vetores inteiros de consumo, e não ordens de revelação. Para $[6,7,7]$, temos $C=(1,2)$. Na camada $m=2$, só existem os consumos abaixo:
\begin{center}
\begin{tabular}{@{}ccccc@{}}
\toprule
$a$ & Cartas restantes & Referências & $|\{j:a_j>0\}|$ & Estados\\
\midrule
$(0,2)$ & $[6]$ & 7 & 1 & $4\cdot1=4$\\
$(1,1)$ & $[7]$ & 6 ou 7 & 2 & $4\cdot2=8$\\
\bottomrule
\end{tabular}
\end{center}
Portanto, $|V_2|=2\cdot2(1+2)=12$. A referência 7 é contada nas duas linhas porque acompanha estoques distintos. Ela pode estar esgotada no estoque: em $(0,2)$, a última carta revelada é 7, mas não resta outro 7 oculto.

Para o final dessa mesma instância, o consumo é $(1,2)$, as referências finais são 6 ou 7 e cada uma admite seis pares $(d,t)$. Assim, $|V_3|=2\cdot3\cdot2=12$. Referência e vez continuam distinguindo terminais, conforme nossa definição de estado. Se $N=0$, há somente a raiz, já terminal; esse caso não utiliza a expressão $2Nk$.

\proxpagina
% PÁGINA 10 --------------------------------------------------------------
\subsection{Folhas e camadas da árvore de históricos}
\begin{proposicao}[Folhas da árvore de históricos]\label{prop:folhas}
Com multiplicidades $C_1,\ldots,C_k$ e $N=\sum_i C_i\ge1$, a árvore de históricos possui
\begin{equation}\label{eq:folhas}
\boxed{\displaystyle L=3^N\frac{N!}{C_1!\cdots C_k!}}
\end{equation}
folhas. Cada folha corresponde a um histórico completo.
\end{proposicao}

\begin{proof}
Um histórico completo é determinado por duas informações: a ordem dos $N$ valores revelados e a sequência das $N$ previsões. Primeiro, se todas as cartas fossem distinguíveis, haveria $N!$ ordens. Permutar entre si as $C_i$ cópias do valor $v_i$ não altera a sequência de valores; cada sequência foi contada $\prod_i C_i!$ vezes. Portanto, existem $N!/\prod_i C_i!$ ordens distintas de valores.

Para cada ordem, há três escolhas de previsão em cada uma das $N$ posições. Como nenhuma previsão é proibida, são $3^N$ sequências. A regra da vez não elimina escolhas: ela apenas determina qual jogador faz cada previsão a partir dos resultados anteriores. O par formado por ordem e previsões determina um histórico único, e todo histórico completo fornece esse par. Trata-se de uma bijeção, de modo que multiplicamos as duas contagens.

Pela Proposição~\ref{prop:representacoes}, um histórico só deixa de admitir continuação depois de consumir todas as cartas. Logo, esses históricos são exatamente as folhas.
\end{proof}

\destaque{Por que o fator é $3^N$ e não $2^N$?}
Na árvore, previsões diferentes continuam distinguindo históricos. Se a próxima carta é 7 e a referência é 6, \emph{menor} e \emph{igual} são duas escolhas diferentes, embora ambas errem e levem ao mesmo estado. As letras $\A/\Er$ resumem os resultados para estudar o grafo; não substituem o registro das três previsões nos históricos.

\destaque{Camadas intermediárias.}
Escrevendo $H_m$ para o número de históricos de comprimento $m$, a mesma bijeção, agora separada por consumo, dá
\begin{equation}\label{eq:historicoscamada}
H_m=3^m\!\sum_{\substack{0\le a_i\le C_i\ (1\le i\le k)\\a_1+\cdots+a_k=m}}
\frac{m!}{a_1!\cdots a_k!},\qquad 0\le m\le N.
\end{equation}
Adotamos $0!=1$, de modo que $H_0=1$. Para cada consumo, o coeficiente multinomial conta suas ordens possíveis; o fator $3^m$ conta as previsões. Em $m=N$, existe apenas $a=C$ e recuperamos~\eqref{eq:folhas}. Por isso, não basta substituir $N$ por $m$ na fórmula das folhas mantendo os $C_i!$ no denominador.

Com $C=(1,2)$, em $m=2$ há as ordens $6,7$; $7,6$; $7,7$. Logo, $H_2=3^2(2+1)=27$. Ao final, são $L=3^3\,3!/(1!2!)=81$ folhas, mas apenas 12 estados terminais. As folhas contam maneiras de a partida acontecer; os terminais contam situações finais distintas.

\proxpagina
% PÁGINA 11 --------------------------------------------------------------
\subsection{Bipartição e classificação dos ciclos}
\begin{teorema}[Bipartição e classificação dos ciclos]\label{teo:ciclos}
O grafo associado ao grafo de estados é bipartido. Ele é uma árvore se, e somente se, $N\le2$. Para $N\ge3$, o menor comprimento de um ciclo simples é 4, qualquer que seja a composição das cartas.
\end{teorema}

\begin{proof}
\textbf{Bipartição.} Separemos $V$ em $\bigcup_{m\text{ par}}V_m$ e $\bigcup_{m\text{ ímpar}}V_m$. Toda aresta liga camadas consecutivas, portanto cruza as duas partes. Um ciclo alterna as partes e só volta à inicial após um número par de arestas. Não há, assim, ciclos ímpares.

\textbf{Até duas cartas.} Para $N=0$, há somente a raiz. Para $N=1$, cada sucessor tem a raiz como único predecessor. Para $N=2$, basta examinar também a camada 2. O estoque final determina as duas cartas consumidas, e a referência final determina qual foi revelada por último. Logo, a ordem de seus valores é recuperada, inclusive quando são iguais. As quatro palavras de resultados têm pares distintos:
\[
\A\A:(2,1),\quad \A\Er:(1,2),\quad
\Er\A:(-1,2),\quad \Er\Er:(0,1).
\]
Assim, o estado final também recupera a palavra de resultados. A primeira carta e o primeiro resultado determinam um único estado na camada 1. As duas previsões erradas já foram reunidas e não criam arestas paralelas. Cada estado não inicial tem, portanto, um único predecessor. O argumento da Proposição~\ref{prop:representacoes}, usando conexidade e pai único, mostra que o associado é uma árvore.

\textbf{Três ou mais cartas.} Fixemos quaisquer três revelações viáveis $x,y,z$, que podem repetir valores. Após $x,y$, as palavras $\Er\Er$ e $\A\Er$ dão estados distintos $u,v$, com mesmo estoque e referência $y$, mas pares $(0,1)$ e $(1,2)$. Revelando $z$, ambos podem chegar a $a$ e $b$, com mesmo novo estoque e referência $z$:
\[
u\xrightarrow{\A}a,\quad u\xrightarrow{\Er}b,\quad
v\xrightarrow{\Er}a,\quad v\xrightarrow{\A}b,
\qquad a:(1,1),\quad b:(0,2).
\]
Os vértices de uma mesma camada diferem em $(d,t)$; os de camadas distintas diferem no total de cartas. Logo, são quatro vértices distintos e, apagadas as orientações, $u,a,v,b,u$ é um ciclo simples de comprimento 4. O grafo é simples e bipartido: não há ciclos simples de comprimentos 1, 2 ou 3. Portanto, esse ciclo é mínimo, e o associado não é uma árvore.
\end{proof}

\begin{center}
\begin{tikzpicture}[>=Stealth,box/.style={draw,rounded corners=2pt,fill=blue!3,minimum width=3.2cm,minimum height=0.85cm,align=center,font=\small}]
\node[box] (u) at (0,1.15) {$u:\ \Er\Er$\\$(d,t)=(0,1)$};
\node[box] (v) at (0,-0.15) {$v:\ \A\Er$\\$(d,t)=(1,2)$};
\node[box] (a) at (6,1.15) {$a:\ \Er\Er\A,\ \A\Er\Er$\\$(d,t)=(1,1)$};
\node[box] (b) at (6,-0.15) {$b:\ \Er\Er\Er,\ \A\Er\A$\\$(d,t)=(0,2)$};
\draw[->] (u)--node[above] {$\A$}(a);
\draw[->] (u)--node[pos=.32,fill=white,inner sep=1pt] {$\Er$}(b);
\draw[->] (v)--node[pos=.32,fill=white,inner sep=1pt] {$\Er$}(a);
\draw[->] (v)--node[below] {$\A$}(b);
\end{tikzpicture}
\captionof{figure}{Camadas 2 e 3. As palavras indicam resultados, com as cartas fixadas. O ciclo existe no associado; nenhuma seta volta à camada anterior.}
\end{center}

\proxpagina
% PÁGINA 12 --------------------------------------------------------------
\subsection{Entradas múltiplas e árvores geradoras}
\begin{proposicao}[Entradas múltiplas e arestas excedentes]\label{prop:excedentes}
Se $o$ é a raiz do grafo de estados, então
\begin{equation}\label{eq:beta}
\boxed{\displaystyle \beta=|E|-|V|+1
=\sum_{v\ne o}\bigl(\grauin(v)-1\bigr).}
\end{equation}
É possível obter uma árvore geradora do grafo associado removendo exatamente $\beta$ arestas. Qualquer remoção que produza uma árvore com todos os vértices precisa retirar essa quantidade de arestas.
\end{proposicao}

\begin{proof}
\textbf{Identidade de contagem.} A raiz não recebe arestas, pois está na camada zero. Todo outro vértice recebe pelo menos uma: ele é alcançado por algum histórico não vazio. Ao somar graus de entrada, cada aresta é contada uma única vez, em seu destino. Portanto,
\[
\sum_{v\ne o}\grauin(v)=|E|,
\qquad
\sum_{v\ne o}(\grauin(v)-1)=|E|-(|V|-1).
\]

\textbf{Construção da árvore geradora.} Para cada $v\ne o$, escolhemos exatamente uma das arestas que chegam a $v$. São $|V|-1$ arestas distintas, pois têm destinos distintos. Sigamos essas arestas escolhidas para trás a partir de qualquer vértice. A camada diminui em um a cada passo, e todo vértice de camada positiva tem uma entrada escolhida. O processo só pode terminar em $o$. Assim, apagadas as orientações, o subgrafo escolhido é conexo e contém todos os vértices.

Ele não tem ciclo simples: um vértice de camada máxima de um suposto ciclo teria duas arestas escolhidas entrando nele, contrariando a seleção de uma única entrada. Logo, o subgrafo é uma árvore geradora. Foram retiradas $|E|-(|V|-1)=\beta$ arestas.

\textbf{Necessidade da quantidade.} Toda árvore finita com $|V|$ vértices tem $|V|-1$ arestas: enraizando-a, cada vértice diferente da raiz corresponde à única aresta que o liga ao pai. Logo, qualquer árvore geradora do nosso grafo mantém exatamente $|V|-1$ arestas e exige a mesma quantidade $\beta$ de remoções.
\end{proof}

\destaque{Interpretação das duas expressões.}
Na fórmula $|E|-|V|+1$, comparamos o número de ligações com o necessário para uma árvore. Na soma dos graus de entrada, reservamos uma entrada para cada estado não inicial e contamos as adicionais. Assim, as entradas múltiplas explicam precisamente as arestas excedentes. Nas notas complementares, as arestas fora de uma árvore geradora são chamadas de \emph{elos}~\cite[p.~130]{rangel}.

Se recolocarmos um desses elos na árvore escolhida, seus extremos já estão ligados por um único caminho simples. A nova aresta fecha um ciclo, chamado fundamental em relação a essa árvore. Desse modo, $\beta$ conta os elos e os ciclos fundamentais dessa construção; não conta todos os ciclos simples, pois outras combinações de arestas podem gerar outros ciclos.

A escolha precisa preservar uma entrada por vértice: retirar $\beta$ arestas arbitrárias pode desconectar o grafo. Além disso, a árvore geradora serve à análise das conexões. Ela deixa de representar algumas jogadas possíveis e, por isso, não substitui o grafo completo para descrever todas as partidas ou resolver decisões do jogo.

\proxpagina
% PÁGINA 13 --------------------------------------------------------------
\section{Exemplos e aplicação das fórmulas}
\subsection{Um final com dois valores distintos}
Considere referência 6, estoque $[6,7,7]$, $d=0$ e $t=1$. Temos $C=(1,2)$, $N=3$ e $k=2$. Pelo Teorema~\ref{teo:vertices},
\[
|V|=1+(3+1)\bigl[1(2+1)+2(1+1)\bigr]=1+4\cdot7=29.
\]
O Corolário~\ref{cor:camadas} dá 12 terminais e a Proposição~\ref{prop:folhas} dá 81 folhas. As camadas permitem acompanhar a diferença entre essas contagens:
\begin{center}
\begin{tabular}{@{}crrr@{}}
\toprule
Jogadas $m$ & Históricos $H_m$ & Estados $|V_m|$ & Arestas de saída\\
\midrule
0 & 1 & 1 & 4\\
1 & 6 & 4 & 12\\
2 & 27 & 12 & 24\\
3 & 81 & 12 & 0\\
\midrule
Total & 115 & 29 & 40\\
\bottomrule
\end{tabular}
\captionof{table}{Árvore e grafo da mesma raiz. Uma folha é um histórico completo.}
\end{center}

Na primeira jogada, os valores 6 e 7 podem ser revelados. Cada um admite acerto e erro: são quatro estados. Na camada 2, os consumos $(0,2)$ e $(1,1)$ admitem, respectivamente, uma e duas referências; cada escolha admite quatro pares $(d,t)$, produzindo $4(1+2)=12$ estados. Na camada 3, o consumo é $(1,2)$ e há duas referências, cada uma com seis pares: são outros 12 estados.

Para contar arestas de saída de uma camada, somamos $2q$ sobre seus estados; $q$ pode variar dentro da camada. A raiz tem $q=2$, logo quatro arestas. Na camada 1, revelar 6 deixa $[7,7]$: seus dois estados têm $q=1$ e contribuem com quatro arestas. Revelar 7 deixa $[6,7]$: seus dois estados têm $q=2$ e contribuem com oito. Na camada 2, cada um dos 12 estados tem uma única carta restante e duas saídas. Assim,
\[
|E|=4+(4+8)+12\cdot2=40,\qquad \beta=40-29+1=12.
\]
As 24 arestas da camada 2 podem chegar a 12 terminais porque origens diferentes podem ter um sucessor comum. Essas arestas continuam distintas, já que uma aresta é determinada por origem e destino.

A árvore contém $1+6+27+81=115$ vértices e, portanto, 114 arestas. O grafo de estados possui 29 vértices e 40 arestas; uma árvore geradora do associado mantém 28 delas. Temos, assim, três objetos distintos: a árvore de históricos, o grafo de estados e uma árvore geradora desse grafo.

\subsection{Três cartas iguais}
Se o estoque passa a ser $[6,6,6]$, ainda com referência 6, temos $k=1$ e $N=3$. Agora $|V|=1+3\cdot4=13$, com camadas de tamanhos $1,2,4,6$. Os sete estados não terminais têm duas saídas cada, logo $|E|=14$ e $\beta=2$. Na árvore, cada histórico não terminal tem três filhos: são $1+3+9+27=40$ vértices e 27 folhas. Mesmo com todas as cartas iguais, previsões que certamente erram permanecem permitidas e produzem diferentes pares de placar e vez.

\proxpagina
% PÁGINA 14 --------------------------------------------------------------
\begin{center}
\begin{tikzpicture}[>=Stealth,x=1cm,y=1cm,
estado/.style={draw,rounded corners=2pt,minimum width=1.15cm,minimum height=.52cm,fill=blue!3,font=\small},
acerto/.style={->,blue!65!black},erro/.style={->,gray!75!black,dashed}]
\node[estado] (o) at (0,3.9) {$(0,1)$};
\node[estado] (a) at (-2.2,2.6) {$(1,1)$};
\node[estado] (e) at (2.2,2.6) {$(0,2)$};
\node[estado] (aa) at (-4,1.3) {$(2,1)$};
\node[estado] (ae) at (-1.4,1.3) {$(1,2)$};
\node[estado] (ee) at (1.4,1.3) {$(0,1)$};
\node[estado] (ea) at (4,1.3) {$(-1,2)$};
\node[estado] (aaa) at (-5,0) {$(3,1)$};
\node[estado] (aae) at (-3,0) {$(2,2)$};
\node[estado] (aee) at (-1,0) {$(1,1)$};
\node[estado] (aea) at (1,0) {$(0,2)$};
\node[estado] (eae) at (3,0) {$(-1,1)$};
\node[estado] (eaa) at (5,0) {$(-2,2)$};
\draw[acerto] (o)--node[above left,font=\scriptsize]{$\A$}(a);
\draw[erro] (o)--node[above right,font=\scriptsize]{$\Er$}(e);
\draw[acerto] (a)--(aa); \draw[erro] (a)--(ae);
\draw[acerto] (e)--(ea); \draw[erro] (e)--(ee);
\draw[acerto] (aa)--(aaa); \draw[erro] (aa)--(aae);
\draw[acerto] (ae)--(aea); \draw[erro] (ae)--(aee);
\draw[acerto] (ee)--(aee); \draw[erro] (ee)--(aea);
\draw[acerto] (ea)--(eaa); \draw[erro] (ea)--(eae);
\foreach \h/\m in {3.9/0,2.6/1,1.3/2,0/3}{\node[anchor=west,font=\small] at (6,\h){$m=\m$};}
\end{tikzpicture}
\captionof{figure}{Grafo completo de $[6,6,6]$ com referência 6. Cada caixa mostra $(d,t)$; a camada determina o estoque e $r=6$ em todos os estados. Linhas contínuas azuis indicam acertos; tracejadas, erros. Há 13 vértices e 14 arestas.}
\end{center}


\subsection{Aplicação ao baralho original de 36 cartas}
A referência inicial está fixada e já foi retirada. Restam $N=35$ cartas: três cópias do valor revelado e quatro de cada um dos outros oito valores. Podemos reordenar as categorias para escrever $C=(3,4,4,4,4,4,4,4,4)$. No Teorema~\ref{teo:vertices},
\begin{align*}
|V|&=1+(35+1)\bigl[3\cdot5^8+8(4\cdot4\cdot5^7)\bigr]\\
&=1+36\cdot11.171.875=402.187.501.
\end{align*}
O fator 36 é $N+1$, e não o número de cartas restantes. A primeira camada tem $2\cdot9=18$ estados e 27 históricos. No final, há
\[
|V_{35}|=2\cdot35\cdot9=630,
\qquad
L=3^{35}\frac{35!}{3!(4!)^8}\approx7{,}83\times10^{44}.
\]
Esses valores vêm das fórmulas, sem enumerar o grafo de centenas de milhões de estados. As folhas contam históricos completos desde a referência fixada; não são resultados finais distintos. Também não são igualmente prováveis por terem sido contadas: as probabilidades dependem do estoque e das escolhas dos jogadores.

\section{Considerações finais}
A passagem da árvore ao grafo reúne passados distintos quando deixam a mesma situação relevante. O invariante do Lema~\ref{lem:pares} controla placar e vez; sua combinação com consumos e referências produz a contagem exata. A diminuição das cartas garante aciclicidade dirigida, enquanto os reencontros entre caminhos produzem ciclos no associado a partir de três cartas. A quantidade $\beta$ relaciona esses reencontros às entradas adicionais e à construção de uma árvore geradora.

Os exemplos mostram como os resultados gerais determinam quantidades de estados, históricos e arestas em composições específicas. A aplicação ao baralho original evidencia o alcance da contagem por fórmula: obtemos o número de vértices sem precisar construir o grafo inteiro. Para composições grandes, mesmo a representação reduzida por estados pode ser impraticável de enumerar.

A visualização dos grafos em HTML complementa a análise teórica e auxilia na compreensão das relações entre os vértices e as arestas. Essa representação se apoia nas definições e nos resultados demonstrados neste relatório.

Comparar estratégias exige acrescentar um objetivo e as probabilidades das revelações. As saídas $\A/\Er$ não são escolhas livres do jogador, que precisa prever antes de ver a carta. Essa extensão pode usar o mesmo estado, mas não é necessária para os resultados estruturais demonstrados neste trabalho.

\begin{thebibliography}{9}
\small
\bibitem{nina} \emph{Notas MD Nina}. Anotações manuscritas de uma aluna da disciplina de Matemática Discreta, ministrada por Maria Soledad Aronna, no ano anterior ao deste trabalho. 57 p. Material de estudo fornecido ao grupo; arquivo local. Autoria identificada apenas pelo nome do arquivo.
\bibitem{rangel} RANGEL, Socorro; OLIVEIRA, Valeriano A. de; ARAUJO, Silvio A. \emph{Elementos de Teoria dos Grafos: Notas de Aula}. Departamento de Matemática Aplicada, IBILCE/Unesp, 24 set. 2018. Disponível em: \url{https://www.ibilce.unesp.br/Home/Departamentos/MatematicaAplicada/docentes/socorro/grafos---notas-de-aula_set2018.pdf}. Acesso em: 13 nov. 2026.
\end{thebibliography}
\end{document}
